% Options for packages loaded elsewhere
\PassOptionsToPackage{unicode}{hyperref}
\PassOptionsToPackage{hyphens}{url}
\PassOptionsToPackage{dvipsnames,svgnames,x11names}{xcolor}
%
\documentclass[
  11pt,
]{article}

\usepackage{amsmath,amssymb}
\usepackage{iftex}
\ifPDFTeX
  \usepackage[T1]{fontenc}
  \usepackage[utf8]{inputenc}
  \usepackage{textcomp} % provide euro and other symbols
\else % if luatex or xetex
  \usepackage{unicode-math}
  \defaultfontfeatures{Scale=MatchLowercase}
  \defaultfontfeatures[\rmfamily]{Ligatures=TeX,Scale=1}
\fi
\usepackage[]{kpfonts}
\ifPDFTeX\else  
    % xetex/luatex font selection
  \setmainfont[]{TeX Gyre Pagella}
\fi
% Use upquote if available, for straight quotes in verbatim environments
\IfFileExists{upquote.sty}{\usepackage{upquote}}{}
\IfFileExists{microtype.sty}{% use microtype if available
  \usepackage[]{microtype}
  \UseMicrotypeSet[protrusion]{basicmath} % disable protrusion for tt fonts
}{}
\makeatletter
\@ifundefined{KOMAClassName}{% if non-KOMA class
  \IfFileExists{parskip.sty}{%
    \usepackage{parskip}
  }{% else
    \setlength{\parindent}{0pt}
    \setlength{\parskip}{6pt plus 2pt minus 1pt}}
}{% if KOMA class
  \KOMAoptions{parskip=half}}
\makeatother
\usepackage{xcolor}
\setlength{\emergencystretch}{3em} % prevent overfull lines
\setcounter{secnumdepth}{-\maxdimen} % remove section numbering


\providecommand{\tightlist}{%
  \setlength{\itemsep}{0pt}\setlength{\parskip}{0pt}}\usepackage{longtable,booktabs,array}
\usepackage{calc} % for calculating minipage widths
% Correct order of tables after \paragraph or \subparagraph
\usepackage{etoolbox}
\makeatletter
\patchcmd\longtable{\par}{\if@noskipsec\mbox{}\fi\par}{}{}
\makeatother
% Allow footnotes in longtable head/foot
\IfFileExists{footnotehyper.sty}{\usepackage{footnotehyper}}{\usepackage{footnote}}
\makesavenoteenv{longtable}
\usepackage{graphicx}
\makeatletter
\def\maxwidth{\ifdim\Gin@nat@width>\linewidth\linewidth\else\Gin@nat@width\fi}
\def\maxheight{\ifdim\Gin@nat@height>\textheight\textheight\else\Gin@nat@height\fi}
\makeatother
% Scale images if necessary, so that they will not overflow the page
% margins by default, and it is still possible to overwrite the defaults
% using explicit options in \includegraphics[width, height, ...]{}
\setkeys{Gin}{width=\maxwidth,height=\maxheight,keepaspectratio}
% Set default figure placement to htbp
\makeatletter
\def\fps@figure{htbp}
\makeatother

\usepackage{luatexja-fontspec}
\IfFontExistsTF{UnBatang.ttf}
  {\newjfontfamily\cvhangul{UnBatang.ttf}}
  {\newjfontfamily\cvhangul{AppleMyungjo}}
\makeatletter
\@ifpackageloaded{caption}{}{\usepackage{caption}}
\AtBeginDocument{%
\ifdefined\contentsname
  \renewcommand*\contentsname{Table of contents}
\else
  \newcommand\contentsname{Table of contents}
\fi
\ifdefined\listfigurename
  \renewcommand*\listfigurename{List of Figures}
\else
  \newcommand\listfigurename{List of Figures}
\fi
\ifdefined\listtablename
  \renewcommand*\listtablename{List of Tables}
\else
  \newcommand\listtablename{List of Tables}
\fi
\ifdefined\figurename
  \renewcommand*\figurename{図}
\else
  \newcommand\figurename{図}
\fi
\ifdefined\tablename
  \renewcommand*\tablename{Table}
\else
  \newcommand\tablename{Table}
\fi
}
\@ifpackageloaded{float}{}{\usepackage{float}}
\floatstyle{ruled}
\@ifundefined{c@chapter}{\newfloat{codelisting}{h}{lop}}{\newfloat{codelisting}{h}{lop}[chapter]}
\floatname{codelisting}{Listing}
\newcommand*\listoflistings{\listof{codelisting}{List of Listings}}
\makeatother
\makeatletter
\makeatother
\makeatletter
\@ifpackageloaded{caption}{}{\usepackage{caption}}
\@ifpackageloaded{subcaption}{}{\usepackage{subcaption}}
\makeatother
\makeatletter
\@ifpackageloaded{fontawesome5}{}{\usepackage{fontawesome5}}
\makeatother
\ifLuaTeX
  \usepackage{selnolig}  % disable illegal ligatures
\fi
\usepackage{bookmark}

\IfFileExists{xurl.sty}{\usepackage{xurl}}{} % add URL line breaks if available
\urlstyle{same} % disable monospaced font for URLs
\hypersetup{
  colorlinks=true,
  linkcolor={blue},
  filecolor={Maroon},
  citecolor={Blue},
  urlcolor={minty},
  pdfcreator={LaTeX via pandoc}}

\usepackage[top=15truemm,bottom=15truemm,left=10truemm,right=10truemm]{geometry}
\usepackage{titlesec}
\usepackage{picture}
\usepackage{fontawesome5}
\usepackage{luatexja}
\definecolor{minty}{HTML}{80c4ac}\definecolor{ParisGreen}{cmyk}{1,0.07,0.10,0.10}

\titleformat{\section}[block]
{}{}{0pt}
{
    \colorbox{minty}{\begin{picture}(0,10)\end{picture}}
    \hspace{0pt}
    \normalfont \Large\bfseries
    \hspace{-4pt}
}
[
\begin{picture}(100,0)
    \put(3,18){\color{minty}\line(1,0){300}}
\end{picture}
\\
\vspace{-30pt}
]

\titlespacing{\section}{0pc}{3.5ex plus .1ex minus .2ex}{1.5ex minus .1ex}

\renewcommand{\labelitemi}{\textcolor{minty}{\faCheckCircle}} %https://mirrors.ibiblio.org/CTAN/fonts/fontawesome/doc/fontawesome.pdf
% \faPaperclip が文献の列挙に良いかも．\author{}
\date{}

\begin{document}

\begin{figure}

\begin{minipage}{0.50\linewidth}
\Huge Hirofumi Shiba\end{minipage}%
%
\begin{minipage}{0.50\linewidth}

\color{minty}

\hfill {\faIcon{home}} \url{https://162348.github.io/}

\par

\hfill {\faIcon{envelope}}
\href{mailto:shiba.hirofumi@ism.ac.jp}{\nolinkurl{shiba.hirofumi@ism.ac.jp}}

\par

\color{gray}\hfill {Last Update: 2026/11/7}\end{minipage}%

\end{figure}%

\vspace{-1em}

\section{Education}\label{education}

\renewcommand{\labelitemi}{\textcolor{minty}{\faGraduationCap}}

\begin{itemize}
\item
  \textbf{PhD Statistical Science}. \emph{Graduate University for
  Advanced Studies, SOKENDAI, Tokyo, Japan}. \hfill {2023.4 -- 2028.3}

  Superivsor:
  Prof.~\href{https://sites.google.com/view/kengokamatani/home}{Kengo
  Kamatani}
\item
  \textbf{BA Mathematics}. \emph{Department of Mathematics, University
  of Tokyo}. \hfill {2019.4 -- 2023.3}

  Supervisor:
  Prof.~\href{https://www.ms.u-tokyo.ac.jp/~nakahiro/hp-naka-e}{Nakahiro
  Yoshida}
\end{itemize}

\section{Publications}\label{publications}

\begin{enumerate}
\def\labelenumi{\arabic{enumi}.}
\tightlist
\item
  \textbf{Hirofumi Shiba}, Kengo Kamatani (2026).
  \href{https://arxiv.org/abs/2602.17087}{Diffusive Scaling Limits of
  Forward Event-Chain Monte Carlo: Provably Efficient Exploration with
  Partial Refreshment}. \emph{Submitted to the Annals of Applied
  Probability}.
\end{enumerate}

\section{Talks}\label{talks}

\begin{enumerate}
\def\labelenumi{\arabic{enumi}.}
\item
  2026-10-13. (15:30-16:30) \textbf{Hirofumi Shiba}.
  \href{https://162348.github.io/posts/2026/Slides/ProbabilityUCL.html}{High-Dimensional
  Behaviour of Two Piecewise Deterministic Markov Processes: Efficiency
  Comparison \& Asymptotic Variance Estimation}.
  \emph{\href{https://sites.google.com/view/probability-at-ucl/hirofumi-shiba}{Bloomsbury
  Probability Seminar}}. 1-19 Torrington Place G08, London, UK.
\item
  2026-06-16. (9:00-10:40) \textbf{Hirofumi Shiba}.
  \href{https://162348.github.io/posts/2026/Slides/IMS-APRM.html}{High-Dimensional
  Behaviour of Piecewise Deterministic Monte Carlo: Efficiency
  Comparison \& Asymptotic Variance Estimation}.
  \emph{\href{https://ims-aprm2026.sta.cuhk.edu.hk/}{The Seventh
  Institute of Mathematical Statistics Asia Pacific Rim Meeting
  (IMS-APRM)}}. Henry Cheng International Conference Centre, Cheng Yu
  Tung Building, The Chinese University of Hong Kong.
\item
  2026-02-10. (9:50-9:55) \textbf{Hirofumi Shiba}.
  \href{https://162348.github.io/posts/2026/Slides/ISI-ISM-ISSAS.html}{Computational
  Complexity Comparison via Scaling Limits for Emerging MCMC Samplers}.
  \emph{\href{https://www3.stat.sinica.edu.tw/3ijc2026/}{ISI-ISM-ISSAS
  Joint Conference}}. Academia Sinica, Taiwan (2nd Conference Room, 3F).
\item
  2026-02-05. (11:40-12:00) 司馬博文.
  \href{https://162348.github.io/posts/2026/Slides/ISM.html}{或る区分確定的マルコフ過程のスケーリング極限}.
  \emph{学生研究発表会（統数研）}. 統数研２階大会議室（B201）.
\item
  2025-07-15. (17:30-18:00) \textbf{Hirofumi Shiba}.
  \href{https://162348.github.io/posts/2025/Slides/DDD.html}{Understanding
  Discrete Denoising Diffusion Models: A Survey and a Gap in the
  Theory}.
  \emph{\href{https://www.uni-ulm.de/en/mawi/mawi-stochastik/forschung/forschungsseminare/forschungsseminar-sommersemester-2025-stochastische-geometrie-und-raeumliche-statistik/}{Workshop
  on Opportunity Week 2025, Stochastische Geometrie und räumliche
  Statistik}}. Universität Ulm.
\item
  2025-02-06. (16:50-17:10) 司馬博文.
  \href{https://162348.github.io/posts/2024/Slides/IRT-ZigZag.html}{動き出す次世代サンプラー:
  区分確定的モンテカルロ}. \emph{学生研究発表会（統数研）}.
  統数研２階大会議室.
\item
  2024-09-10. \textbf{Hirofumi Shiba} \textbar{} 司馬博文.
  \href{https://162348.github.io/posts/2024/Slides/ZigZagPoliSci.html}{Zig-Zag
  Sampler: A MCMC Game-Changer}. \emph{Talk}. Seoul National University,
  Gwanak ({\cvhangul 관악}) campus, South Korea.
\item
  2024-07-25. 司馬博文.
  \href{https://162348.github.io/posts/2024/Slides/ZigZagSampler.html}{Zig-Zag
  サンプラー: 物理のくびきを超える MCMC}. \emph{ベイズ会}.
  本郷キャンパス小島ホール第二セミナー室.
\end{enumerate}

\section{Posters}\label{posters}

\begin{enumerate}
\def\labelenumi{\arabic{enumi}.}
\item
  2026-10-31. (14:00-15:30) 司馬博文.
  \href{https://162348.github.io/posts/2026/Posters/\%E7\%95\%B0\%E5\%88\%86\%E9\%87\%8E\%E7\%95\%B0\%E6\%A5\%AD\%E7\%A8\%AE\%E7\%A0\%94\%E7\%A9\%B6\%E4\%BA\%A4\%E6\%B5\%81\%E4\%BC\%9A.html}{Sticky
  Manifold PDMP による低次元多様体上に台を持つ分布からのサンプリング}.
  \emph{\href{https://jsiam.org/jsiam_archive/kr/career2026/}{異分野異業種研究交流会}}.
  東京大学 駒場キャンパスⅠ 21 KOMCEE West 地下1階MMホール.
\item
  2026-07-02. (13:50-15:20) \textbf{Hirofumi Shiba}.
  \href{https://162348.github.io/posts/2026/Posters/ISBA_Satellite.html}{A
  PDMP Scaling Analysis: Gains from Persistent Momentum}.
  \emph{\href{https://kkamatani.github.io/2026-isba-satellite/}{ISBA
  Satellite Meeting 2026}}. Institute of Statistical Mathematics.
\item
  2026-06-29. (17:15-19:30) \textbf{Hirofumi Shiba}.
  \href{https://162348.github.io/posts/2026/Posters/ISBA.html}{A
  Provable Comparison Scheme via Scaling Limits for Emerging Monte Carlo
  Samplers}. \emph{\href{https://isba2026.github.io/}{ISBA World Meeting
  2026}}. WINC AICHI (Aichi Industry \& Labor Center).
\item
  2026-05-22. (13:00-15:00) 司馬博文.
  \href{https://162348.github.io/posts/2026/Posters/OH.html}{PDMP
  の拡散極限と早期収束診断}.
  \emph{\href{https://www.ism.ac.jp/openhouse/2026/index.html}{統計数理研究所オープンハウス}}.
  １階交流アトリウム.
\item
  2026-03-03. (16:00-18:00) \textbf{Hirofumi Shiba}.
  \href{https://162348.github.io/posts/2026/Posters/FIMI.html}{Diffusive
  Scaling Limits \& Early Diagnostics for PDMP Samplers}.
  \emph{\href{https://fimiws.github.io/fimi2026/}{Workshop on Functional
  Inference and Machine Intelligence}}. TKP Shinjuku West Exit
  Conference Center, Shinjuku, Tokyo, Japan.
\item
  2026-02-16. (15:45-16:45) 司馬博文.
  \href{https://162348.github.io/posts/2026/Posters/ISACT.html}{区分確定的モンテカルロ法の拡散極限と早期収束診断}.
  \emph{\href{https://okuno.net/events/ISACT2026/}{第２回：計算技術による学際的統計解析ワークショップ}}.
  Institute of Statistical Mathematics (Room D305).
\item
  2026-02-10. (8:50--11:00, 11:20-12:20 (Group B)) \textbf{Hirofumi
  Shiba}.
  \href{https://162348.github.io/posts/2026/Posters/ISI-ISM-ISSAS.html}{Computational
  Complexity Comparison via Scaling Limits for Emerging MCMC Samplers}.
  \emph{\href{https://www3.stat.sinica.edu.tw/3ijc2026/}{ISI-ISM-ISSAS
  Joint Conference}}. Academia Sinica, Taiwan (2nd Conference Room, 3F).
\item
  2025-11-12. (17:10-19:10 (1b-R-02)) 司馬博文.
  \href{https://162348.github.io/posts/2025/Posters/IBIS.html}{メトロポリスを超えた枠組みで我々はどこまで行けるか？:
  Forward Event-Chain Monte Carlo のスケーリング解析}.
  \emph{\href{https://ibisml.org/ibis2025/posters/}{IBIS 2025}}.
  那覇文化芸術劇場なはーと.
\item
  2025-09-22. (15:00-16:30) 司馬博文.
  \href{https://162348.github.io/posts/2025/Posters/FECMC.html}{メトロポリスを超えた枠組みで我々はどこまで行けるか？:
  PDMP と従来の MCMC はどう違い，どう違わないか？}.
  \emph{\href{https://sites.google.com/view/yaml2025/}{機械学習若手の会
  (YAML) 2025}}. ハートピア熱海.
\item
  2025-08-28. (17:20-19:00) 司馬博文.
  \href{https://162348.github.io/posts/2025/Posters/Acceleration.html}{サンプリングにおける加速と安定性}.
  \emph{\href{https://www.ism.ac.jp/~mirai/sscoke/2025/}{連続最適化および関連分野に関する夏季学校}}.
\item
  2025-05-23. (13:40-14:20) 司馬博文.
  \href{https://162348.github.io/posts/2025/Posters/OH.html}{ベイズ変数選択の計算的解決:
  PDMP による非絶対連続分布からのサンプリング}.
  \emph{\href{https://www.ism.ac.jp/openhouse/}{統計数理研究所オープンハウス}}.
\item
  2025-03-08. (12:00-12:50) 司馬博文.
  \href{https://162348.github.io/posts/2025/Posters/NewRegime.html}{PDMP
  によりスパイク付きの非絶対連続分布からもサンプリングが可能}.
  \emph{\href{https://jss2025spring.ywstat.jp/}{第19回日本統計学会春季集会}}.
  (優秀発表賞受賞)
\item
  2025-02-17. (15:30-16:30) 司馬博文.
  \href{https://162348.github.io/posts/2025/Posters/Sticky.html}{ベイズ変数選択の計算的解決:
  PDMP による非可逆ジャンプの達成}.
  \emph{\href{https://okuno.net/events/ISACT2025\#poster-speakers}{計算技術による学際的統計解析ワークショップ}}.
\item
  2024-05-24. (10:30-11:10) 司馬博文.
  \href{https://162348.github.io/posts/2024/Posters/OH.html}{新時代の
  MCMC を迎えるために: 連続時間アルゴリズムへの進化}.
  \emph{\href{https://www.ism.ac.jp/openhouse/2024/index.html}{統数研オープンハウス}}.
\item
  2024-02-25. (18:00-19:30) \textbf{Hirofumi Shiba}.
  \href{https://162348.github.io/posts/2024/Posters/RecentDevelopment.html}{A
  Recent Development of Particle Methods: Inquiry towards a Continuous
  Time Limit and Scalability}.
  \emph{\href{https://groups.oist.jp/mlss}{MLSS2024 (OIST, Okinawa,
  Japan)}}.
\end{enumerate}

\section{Research Visit}\label{research-visit}

\renewcommand{\labelitemi}{\textcolor{minty}{\faUniversity}}

\begin{itemize}
\item
  School of Mathematics, \textbf{University of Bristol}.
  \hfill {2026.9.21 -- 2026.10.21}

  Supervisor: Prof.~Christophe Andrieu
\item
  Institute of Stochastics, \textbf{Ulm University}, Germany.
  \hfill {2025.7.13 -- 2025.7.26}

  Supervisor: Prof.~Evgeny Spodarev
\item
  Department of Statistics \& Data Science, \textbf{National University
  of Singapore}. \hfill {2025.6.1 -- 2025.6.30}

  Supervisor: Dr.~Alexandre Thiéry
\item
  Department of Statistical Science, \textbf{University College London},
  United Kingdom. \hfill {2024.11.4 -- 2024.12.2}

  Supervisor: Prof.~Alexandros Beskos
\end{itemize}

\section{Grants \& Awards}\label{grants-awards}

\renewcommand{\labelitemi}{\textcolor{minty}{\faAward}}

\begin{itemize}
\item
  \textbf{Research Fellow}. \emph{Graduate University for Advanced
  Studies, SOKENDAI, Tokyo, Japan}. \hfill {2025.4 -- 2028.3}

  Supported by JST BOOST, Japan Grant Number JPMJBS2412: \emph{A Radical
  Approach to Trustworthy AI in the Coming Era}. ¥ 11,700,000.
\item
  \textbf{Best Poster Award}, 19th Spring Meeting of the Japan
  Statistical Society. \hfill {2025.3.8}
\item
  \textbf{Director's Award} from the Japan Statistical Society
  Certificate. \hfill {2025.3.8}
\end{itemize}

\renewcommand{\labelitemi}{\textcolor{minty}{\faBookmark}}

\section{Service}\label{service}

\textbf{Reviewing}: Bayesian Analysis, Journal of Machine Learning
Research

\section{Additional Experience}\label{additional-experience}

\renewcommand{\labelitemi}{\textcolor{minty}{\faSuitcase}}

\begin{itemize}
\item
  \textbf{Data Scientist}. \emph{PreMedica, Inc., Tokyo, Japan}.
  \hfill {2024.9 -- today}

  \emph{Responsibilities}: Provided Bayesian data analysis solutions to
  healthcare applications.
\item
  \textbf{Research Assistant}. \emph{The Institute of Statistical
  Mathematics, Tokyo, Japan}. \hfill {2023.7 -- 2026.3}

  \emph{Responsibilities}: Contributed to the R package \texttt{YUIMA},
  an open-source project aiming to simulate and infer multidimensional
  stochastic differential equations, with an emphasis on Bayesian
  inference.
\end{itemize}



\end{document}
